\chapter[Nachrichten aus unserer Indianermission]{Nachrichten aus unserer Indianermission in Wisconsin}
\quellseite{4}

Was die ältere Mission unter den Stockbridge-Indianern betrifft, so macht das Werk des Herrn unter diesen zwar langsame, aber doch sichtliche Fortschritte. Der Kirchenbesuch kann wohl ein guter genannt werden. Missionar Kretzmann berichtete kürzlich, daß durchschnittlich 50 bis 55 Personen die Gottesdienste besucht hätten. Auch der Schulbesuch nimmt sichtlich zu. Die Schule wurde zuletzt von 29 Kindern besucht, und 5 werden noch in der nächsten Zeit erwartet. Mit der Hauptgemeinde ist noch eine Filiale verbunden, sieben Meilen entfernt, in dem sogenannten Hinterbuch. Der Missionar geht alle vier Wochen des Nachmittags dorthin. Der Besuch des Gottesdienstes dort war Anfangs nicht der beste, doch geht es nach und nach etwas besser. Das letzte Mal waren 20 Personen zugegen.

Die Missionsfreunde werden aber gewiß gern auch etwas von dem \unsicher{Woriaskett-} oder Heidensettlement, wie es auf der Reservation genannt wird, hören wollen. Dieses Settlement wird mit vollem Rechte das \gesperrt{Heiden}settlement genannt; denn wie die Heiden in ihren Ländern groben Götzendienst treiben, so geschieht das auch noch in der Nordwestecke von Shawano County, Wis., wo noch diesen Herbst ein Götzenfest und Götzentänze stattgefunden haben. Wir haben also auch in unserm
\quellseite{5}
eigenen Lande Gelegenheit und Veranlassung, eigentliche Heidenmission zu treiben.

Daß die Schule und die Wohnung für einen Lehrer in jenem Settlement fertiggestellt sind, werden die Missionsfreunde wohl schon erfahren haben. Wir hatten dann auch zu unserer Freude einen Mann gefunden, von dem wir hofften, daß unter ihm die Schule einen guten Anfang nehmen werde. Der Anfang wurde auch gemacht; aber leider war der Lehrer nicht der rechte Mann; denn als es galt, einige Schwierigkeiten zu überwinden und Unannehmlichkeiten zu tragen, verließ er ganz plötzlich die Station. Missionar Kretzmann und auch unsere Missionscommission kam dadurch in große Verlegenheit. Die Schule hatte ihren Anfang genommen und mußte nothgedrungen fortgesetzt werden. Der Indianeragent hatte verlangt, daß unsere Schule am 1. September 1903 beginne, sonst werde er alle Kinder von dort nach St. in die Regierungsaufschule bringen lassen. Sollte diese Drohung nicht ausgeführt werden, so durfte unsere Schule nicht unterbrochen werden. In unserer Noth wußten wir uns nicht anders zu helfen, als daß wir auf Rath Missionar Kretzmanns eine christliche Indianerfrau vorläufig als Lehrerin anstellten. Diese Frau ist seiner Zeit von P. Nickel getauft und später von Missionar Kretzmann confirmirt worden. Sie besitzt in den weltlichen Fächern ziemliche Kenntnisse, hat auch einigermaßen eine christliche Erkenntniß und führt auch einen christlichen Wandel. Sie unterrichtet die Kinder, 11 an der Zahl, in den weltlichen Fächern und erzählt ihnen biblische Geschichten. So steht es jetzt noch mit der Schule dort; denn so viel Mühe wir uns auch seither gegeben haben, einen Lehrer dafür zu gewinnen, so ist es doch bisher alles vergeblich gewesen. Alle Berufe wurden abgelehnt. Vielleicht denkt man sich diesen Beruf auch noch schwerer, als er ist.

Unser Herr Christus fordert von allen seinen Jüngern, seinen Christen, Selbstverleugnung. Es ist auch gewiß wahr, daß dort von einem Lehrer oder Missionar viel Selbstverleugnung geübt werden muß und auch mancherlei Entbehrungen zu tragen sind. Der nächste Platz, wo sich das Postamt befindet und wo Weiße wohnen, ist sechs Meilen, die nächste Kirche unseres Glaubens vierzehn Meilen entfernt, das heißt, auf den bis jetzt vorhandenen Wegen. Aber wahrer Christen Art ist, daß sie nicht sehen auf das Sichtbare, sondern auf das Unsichtbare. Ihnen gilt auch der Umgang und die Gemeinschaft mit Gott viel höher als der Umgang und die Gemeinschaft mit Menschen. Es ist gewiß und steht uns Christen außer allem Zweifel: wenn diesen Heiden nicht das Evangelium von Christo gebracht wird, dann fahren sie ins ewige Verderben. Jedermann, der dies durch die Gnade Gottes verhindern und abwenden kann, ist dazu verbunden. Und welch selige Freude ist es, wenn man die Wahrheit des Wortes erfahren darf: „Wer den Sünder bekehret hat von dem Irrthum seines Weges, der hat einer Seele vom Tode geholfen“, Jac. 5, 20. Diese selige Freude wiegt gewiß hernach alle Opfer der Selbstverleugnung nicht nur auf, sondern geht weit darüber hinaus. Und auch im Irdischen wird es dort bald etwas besser werden. Schon jetzt kann der Lehrer oder Missionar von Zeit zu Zeit Conferenzen mit Amts- und Glaubensbrüdern besuchen, wo er sich Rath und Trost holen und brüderliche Gemeinschaft pflegen kann. Er ist auch nach Süden, Osten und Westen zu in keiner großen Entfernung von vielen Amts- und Glaubensbrüdern umgeben.

Schließlich möchte noch der Unterzeichnete im Auftrage der Missionscommission an alle Leser dieser Zeilen die herzliche Bitte richten: Wer uns einen geeigneten Lehrer oder Missionar zur Berufung für diese Mission nennen kann, der wolle dieses doch unter folgender Adresse thun: Rev. H. Erck, 150 Bowen St., Oshkosh, Wis.

Im Auftrage der Indianermissionscommission
P. H. Dicke.
